\documentclass[10pt, aspectratio=169]{beamer}
\input{../../_en_preambulo_beamer}
% Comandos y macros estadísticas compartidas para las presentaciones Beamer en inglés (Metropolis)

% Conjuntos numéricos básicos
\newcommand{\R}{\mathbb{R}}
\newcommand{\N}{\mathbb{N}}
\newcommand{\Q}{\mathbb{Q}}
\newcommand{\Z}{\mathbb{Z}}
\newcommand{\set}[1]{\left\{ #1 \right\}}
\newcommand{\abs}[1]{\left|#1\right|}
\newcommand{\norm}[1]{\left\| #1 \right\|}

% Comandos estadísticos y probabilísticos del curso
\newcommand{\Var}{\operatorname{Var}}
\newcommand{\cov}{\operatorname{Cov}}
\newcommand{\corr}{\rho}
\newcommand{\comb}[2]{\begin{pmatrix} #1 \\ #2 \end{pmatrix}}
\newcommand{\s}{\sigma}
\newcommand{\particion}{\mathfrak{P}}
\newcommand{\card}[1]{\#\left( #1 \right)}
\newcommand{\seq}[3]{\set{#1}_{#2}^{#3}}
\newcommand{\E}{\mathbb{E}}
\newcommand{\Prob}{\mathbb{P}}

% Entornos pedagógicos y cajas en inglés académico natural (soportando tanto nombres en inglés como en español)
\newenvironment{discussion}{%
  \begin{alertblock}{Discussion Question}%
}{%
  \end{alertblock}%
}
\newenvironment{discusion}{%
  \begin{alertblock}{Discussion Question}%
}{%
  \end{alertblock}%
}

\newenvironment{reflection}{%
  \begin{alertblock}{Classroom Reflection}%
}{%
  \end{alertblock}%
}
\newenvironment{reflexion}{%
  \begin{alertblock}{Classroom Reflection}%
}{%
  \end{alertblock}%
}

\newenvironment{paradox}{%
  \begin{alertblock}{Exploratory Paradox}%
}{%
  \end{alertblock}%
}
\newenvironment{paradoja}{%
  \begin{alertblock}{Exploratory Paradox}%
}{%
  \end{alertblock}%
}

\newenvironment{motivation}{%
  \begin{exampleblock}{Motivation: Data Science in Practice}%
}{%
  \end{exampleblock}%
}
\newenvironment{motivacion}{%
  \begin{exampleblock}{Motivation: Data Science in Practice}%
}{%
  \end{exampleblock}%
}

\newenvironment{quickchallenge}{%
  \begin{block}{Quick Team Challenge}%
}{%
  \end{block}%
}
\newenvironment{retoclase}{%
  \begin{block}{Quick Team Challenge}%
}{%
  \end{block}%
}


\title[Optimal Coefficients]{Section 09.05: Optimal Coefficients and Significance Tests}
\subtitle{The $t$-Test for the Slope, the $F=t^2$ Identity, and the Residual Standard Error}
\author[J. Castillo Colmenares]{Juliho Castillo Colmenares}
\institute{Tecnológico de Monterrey}
\date{\vspace{-1.5cm}}

\begin{document}

% Slide 1: Title
\begin{frame}[plain]
  \titlepage
\end{frame}

% Slide 2: Roadmap
\begin{frame}{Roadmap --- Chapter 09: Linear and Multiple Regressions}
  \scriptsize
  Where are we in the modular development of this chapter?
  \begin{itemize}\itemsep=0.02em
    \item \textbf{09.01} Correlation as the Premise of Regression
    \item \textbf{09.02} Introduction to Linear Regression
    \item \textbf{09.03} The Mathematics of Regression: OLS Derivation and $R^2$
    \item \textbf{09.04} Linear Regression on Simulated Data
    \item \textbf{09.05} Optimal Coefficients, $t$/$F$ Tests, and RSE \emph{(Today)}
    \item \textbf{09.06} Implementation in Python with \texttt{statsmodels}
    \item \textbf{09.07} Multiple Regression, the Normal Equation, and Ridge/Lasso
    \item \textbf{09.08} Model Validation and $k$-fold Cross-Validation
    \item \textbf{09.09} Residual Diagnostics and Multicollinearity (VIF)
    \item \textbf{09.10} Regression with \texttt{scikit-learn}
    \item \textbf{09.11} Categorical Variables and Dummy Variables
    \item \textbf{09.12} Nonlinear Transformations and Polynomial Regression
  \end{itemize}
  \pause
  \vspace{-0.05cm}
  \begin{block}{Session Objective}
    Subject the estimated slope coefficient to a formal hypothesis test, verify the algebraic identity $F=t^2$ connecting individual and global significance in simple regression, and compute the Residual Standard Error (RSE) as an absolute --- not relative --- measure of fit quality.
  \end{block}
\end{frame}

% Slide 3: Motivation
\begin{frame}{Are the Estimated Coefficients Trustworthy?}
  \footnotesize
  \begin{motivacion}
    When we compute the values of $\hat\alpha$ and $\hat\beta$, we obtain \textbf{estimates}, not exact values: each carries an associated $p$-value. We need to demonstrate their statistical significance via a hypothesis test before trusting the model.
  \end{motivacion}

  \vspace{0.15cm}
  \pause
  The central question: is the $\hat\beta$ we obtained different from zero, or could it simply be sample noise around $\beta=0$ (no real linear relationship)?
\end{frame}

% Slide 4: Hypothesis test for beta
\begin{frame}{Hypothesis Test for the Slope}
  \footnotesize
  \begin{block}{Setup}
    $$H_0:\beta=0 \quad (\text{no linear relationship}) \qquad \text{vs.} \qquad H_a:\beta\neq0$$
  \end{block}
  \pause

  \vspace{0.1cm}
  \begin{alertblock}{Test Statistic}
    $$t=\frac{\hat\beta}{SE(\hat\beta)}, \qquad SE(\hat\beta)=\sqrt{\frac{\text{CME}}{S_{XX}}} \ \sim\ t_{n-2} \quad\text{under } H_0$$
  \end{alertblock}
\end{frame}

% Slide 5: The F=t² identity
\begin{frame}{The $F=t^2$ Identity in Simple Regression}
  \footnotesize
  \begin{block}{Global $F$-Test}
    $$F=\frac{\text{SSR}/1}{\text{SSE}/(n-2)} \sim F_{1,n-2}$$
  \end{block}
  \pause

  \vspace{0.1cm}
  \begin{alertblock}{Exact Connection to the $t$-Test}
    In simple regression ($p=1$ predictor), the global $F$-test for model significance and the individual $t$-test for the slope are \textbf{mathematically equivalent}: $F=t^2$. With multiple predictors (Section 09.07), this equivalence disappears: the $F$ statistic tests \emph{joint} significance, while each $t_j$ tests \emph{individual} significance.
  \end{alertblock}
\end{frame}

% Slide 6: Residual Standard Error
\begin{frame}{The Residual Standard Error (RSE)}
  \footnotesize
  \begin{block}{Definition}
    $$RSE=\sqrt{\frac{\text{SSE}}{n-2}}$$
  \end{block}
  \pause

  \vspace{0.1cm}
  \begin{alertblock}{Interpretation}
    RSE estimates the standard deviation of the error term: it is the typical prediction error \emph{in the original units of $Y$}, even if the coefficients were known with absolute precision. It is interpreted by comparing it against $\bar Y$: $RSE/\bar Y$ gives a percentage relative error.
  \end{alertblock}
\end{frame}

% Slide 7: Python Lab (Block 1)
\begin{frame}[fragile]{Python Lab: Optimal Coefficients and the $t$-Test}
  \scriptsize
  OLS fit and slope significance test (\texttt{09.05\_optimal\_coefficients\_and\_tests.py}):
  {\lstset{basicstyle=\fontsize{3.4pt}{4.1pt}\selectfont\ttfamily,aboveskip=0.05em,belowskip=0.05em}
  \lstinputlisting[firstline=18, lastline=35]{../../code/09_regresiones/09.05_optimal_coefficients_and_tests.py}}
\end{frame}

% Slide 8: Python Lab (Block 2)
\begin{frame}[fragile]{Python Lab: Verifying $F=t^2$}
  \scriptsize
  Numeric confirmation of the exact algebraic identity between the $t$- and $F$-tests:
  {\lstset{basicstyle=\fontsize{3.6pt}{4.3pt}\selectfont\ttfamily,aboveskip=0.05em,belowskip=0.05em}
  \lstinputlisting[firstline=41, lastline=51]{../../code/09_regresiones/09.05_optimal_coefficients_and_tests.py}}
\end{frame}

% Slide 9: Python Lab (Block 3)
\begin{frame}[fragile]{Python Lab: Residual Standard Error}
  \scriptsize
  Computing the RSE and its interpretation as a relative error:
  {\lstset{basicstyle=\fontsize{3.8pt}{4.5pt}\selectfont\ttfamily,aboveskip=0.05em,belowskip=0.05em}
  \lstinputlisting[firstline=54, lastline=63]{../../code/09_regresiones/09.05_optimal_coefficients_and_tests.py}}
\end{frame}

% Slide 10: Terminal Output
\begin{frame}[fragile]{Python Lab: Terminal Output}
  \scriptsize
  Standard output produced when running the lab script:
  \vspace{0.02cm}
  \begin{block}{Standard Console Output}
    \fontsize{3.6pt}{4.3pt}\selectfont
    \begin{verbatim}
=== Block 1: Optimal Coefficients and t-Test ===
alpha_hat=1.8577, beta_hat=0.3468
SE(beta_hat)=0.0202, t=17.1624 (df=98), p-value=2.73e-31

=== Block 2: F-Test (F = t^2) ===
F statistic = 294.5482 (df=1,98), p-value=2.73e-31
t^2 = 294.5482, |F - t^2| = 0.0000000000

=== Block 3: Residual Standard Error ===
RSE = 0.4453
mean(Y) = 2.4203
Relative error = RSE/mean(Y) = 0.1840 (18.4%)
    \end{verbatim}
  \end{block}
\end{frame}


% Slide 13: Closing
\begin{frame}{Section Closing: Optimal Coefficients and Significance Tests}
  \begin{reflexion}
    An estimated coefficient without a significance test is just a number: we have formalized when $\hat\beta$ represents real evidence of a linear relationship, discovered the elegant $F=t^2$ identity unique to simple regression, and established the RSE as the absolute reference measure for judging any model's predictive quality.
  \end{reflexion}

  \vspace{0.2cm}
  \pause
  \begin{block}{Key Takeaways}
    \begin{itemize}\itemsep=0.08cm
      \item \textbf{$t$-test:} $t=\hat\beta/SE(\hat\beta)\sim t_{n-2}$ under $H_0:\beta=0$.
      \item \textbf{$F=t^2$:} exact equivalence exclusive to single-predictor regression.
      \item \textbf{RSE:} absolute prediction error in $Y$'s units; $RSE/\bar Y$ gives the relative error.
    \end{itemize}
  \end{block}
\end{frame}

% Slide 14: Modular Perspective
\begin{frame}{Modular Perspective --- The Next Challenge}
  \scriptsize
  \begin{retoclase}
    We have computed everything by hand with closed-form formulas. Section 09.06 shows how to obtain exactly the same results --- coefficients, $p$-values, $R^2$, RSE --- with a single line of professional code using the \texttt{statsmodels} library, the industry-standard tool for statistical regression in Python.
  \end{retoclase}

  \vspace{0.08cm}
  \pause
  \begin{alertblock}{Recommended Preparation}
    \begin{itemize}\itemsep=0.06cm
      \item Install \texttt{statsmodels} if you don't have it yet (\texttt{pip install statsmodels}).
      \item Reflect on the advantages and risks of using a "black-box" library vs. implementing the formulas manually as in this section.
    \end{itemize}
  \end{alertblock}
\end{frame}

\end{document}
